\cotefile{137}
\scanpage{2}
$k$ anneau commutatif, $k_*$ \unsure{le $k$ simplicial|le demi-groupe|est l'anneau $\frac12$ simpl.} \unsure{défini}.

\uline{Lemme 1} Soit $L_*$ un $k$-module $\frac12$ simplicial, \unsure{plat et f.g.}

Pour que $L_*$ soit un $k_*$-Module projectif, il suffit qu'il soit flasque (i.e. $\pi_i(L_*)=0$ $\forall i\geq0$) et que les $L_n$ soient des $k$-modules projectifs.
% margin note written sideways
\margin{\unsure{p.ex. il} suffit \unsure{que les} $\unsure{L_q(L_*)}$ \unsure{soient plats}.}

La \unsure{condition|complexe} signifie aussi que le complexe des chaînes (des $k$-modules) associé par Dold-Puppe est acyclique et à \unsure{compos.} projectives, on voit alors qu'il \deleted{est} \inserted{\unsure{devient}} isomorphe à une somme directe de complexes de chaînes de la forme
\[
\ldots\,0\longrightarrow0\longrightarrow L_q
 \xrightarrow{\unsure{\sim}}L_q\longrightarrow0\longrightarrow0\,\ldots
\]
\unsure{non dégénérés|non dégénéres} \unsure{à fois conséc.|en degrés conséc.}, en $\unsure{k}$. Or un tel complexe est évidemment projectif (dans la cat. des complexes de chaînes de $k$-modules), \unsure{donc aussi une somme quelconque de tels}, d'où on conclut \unsure{par l'équivalence de catégories}.

\uline{NB} Il suffit, au lieu de supposer $L_*$ flasque i.e. $\pi_i(L_*)=0$ $\forall i\geq0$, de supposer $\pi_i(L_*)=0$ pour $i>0$ et $\pi_0(L_*)$ $k$-projectif. (On trouve \unsure{des sommes directes des complexes|alors sans doute une condition} \unsure{en} \unsure{ci et s.}) En effet, dans l'argument précédent, il faut de plus tenir compte \unsure{d'un sommande|des sommandes} qui est un module projectif $L_0$ placé en degré $0$. Or il est immédiat que c'est encore un complexe de chaînes projectif.

\uline{Cor.} Les $\Gamma^i\Phi_{*,k}\otimes_k\Lambda^j\Psi_{*,k}$ \inserted{et $\Lambda^i\Phi_{*,k}\otimes\Lambda^j\Psi_{*,k}$} (pour $i>0$ sont des modules $\frac12$ simpliciaux projectifs.
% archive sheet number: 1
